\chapter{uv の仕組みを理解する}

\section{uv の概要}

\textbf{uv}（読み方：ユーブイ）は Astral 社が開発した
Python パッケージマネージャです。
公式サイト：\url{https://docs.astral.sh/uv/}

従来の \cmd{pip}（パッケージインストーラー）と
\cmd{venv}（仮想環境作成ツール）の両方の機能を一つに統合しており、
以下の特徴があります。

\begin{description}
  \item[\textbf{高速}]
    Rust で書かれており、\cmd{pip} より 10〜100 倍速い。
    多数のパッケージをインストールしても数秒で完了する
  \item[\textbf{再現性が高い}]
    \cmd{uv.lock} ファイルで全員がまったく同じバージョンを使える
  \item[\textbf{Python も自動管理}]
    Python 本体のインストールも自動で行う。
    事前に Python をインストールしておく必要がない
  \item[\textbf{シンプル}]
    \cmd{uv sync} 一発で仮想環境の作成からパッケージのインストールまで完了する
\end{description}

\section{pip + venv との比較}

従来の方法（\cmd{pip} + \cmd{venv}）と \cmd{uv} を比べると次のとおりです。

\begin{center}
\begin{tabularx}{\linewidth}{lXX}
  \toprule
  操作 & pip + venv & uv \\
  \midrule
  Python のインストール & 別途必要 & 自動 \\
  仮想環境を作る & \cmd{python3 -m venv .venv} & 自動（\cmd{uv sync} で） \\
  仮想環境を有効化 & \cmd{source .venv/bin/activate} & 不要 \\
  パッケージをインストール & \cmd{pip install -r req.txt} & \cmd{uv sync} \\
  スクリプトを実行 & \cmd{python script.py} & \cmd{uv run python script.py} \\
  パッケージを追加 & \cmd{pip install xxx} & \cmd{uv add xxx} \\
  新規プロジェクト作成 & 手動で \cmd{pyproject.toml} 作成 & \cmd{uv init} \\
  速度 & 普通 & 10〜100 倍速い \\
  \bottomrule
\end{tabularx}
\end{center}

\begin{tipbox}[activate が不要な理由]
uv では \cmd{source .venv/bin/activate} による仮想環境の有効化が不要です。
\cmd{uv run} をつけるだけで、自動的に \cmd{.venv/} 内の Python と
ライブラリを使って実行してくれます。
「有効化し忘れてシステム Python で実行してしまった」という
よくあるミスがなくなります。
\end{tipbox}

% -------------------------------------------------
\section{pyproject.toml とは}
% -------------------------------------------------

\cmd{pyproject.toml} はプロジェクトの設定ファイルです。
「このプロジェクトには何が必要か」をまとめて書く場所です。

\subsection{ファイルの構造}

B4課題フォルダの \cmd{pyproject.toml} を見てみましょう。

\begin{lstlisting}[style=toml, caption={pyproject.toml の構造}]
[project]
name = "nmrlab-b4-2026"      # プロジェクトの名前
version = "2026.1"            # バージョン
requires-python = ">=3.11"   # Python 3.11 以上が必要

dependencies = [             # 必要なライブラリの一覧
    "numpy>=1.26",           # numpy バージョン 1.26 以上
    "scipy>=1.12",
    "matplotlib>=3.8",
    "scikit-image>=0.22",
    "PyWavelets>=1.5",
    "bm3d>=4.0",
    "openpyxl>=3.1",
    "jupyterlab>=4.0",
    "ipywidgets>=8.0",
    "japanize-matplotlib>=1.1",
]

[tool.uv]
dev-dependencies = [         # 開発時だけ使うツール（通常不要）
    "pytest>=8.0",
]
\end{lstlisting}

\subsection{バージョン指定の書き方}

\begin{center}
\begin{tabularx}{\linewidth}{lX}
  \toprule
  書き方 & 意味 \\
  \midrule
  \cmd{"numpy>=1.26"} & 1.26 以上ならどのバージョンでも OK \\
  \cmd{"numpy==1.26.4"} & 1.26.4 ぴったり（固定） \\
  \cmd{"numpy>=1.26,<2.0"} & 1.26 以上かつ 2.0 未満 \\
  \cmd{"numpy"} & バージョン指定なし（最新） \\
  \bottomrule
\end{tabularx}
\end{center}

\begin{notebox}[pyproject.toml は「希望」を書く場所]
\cmd{pyproject.toml} の \cmd{dependencies} は
「このバージョン以上なら動く」という\textbf{条件}を書くところです。
実際にインストールする\textbf{正確なバージョン}は
\cmd{uv.lock} に記録されます。
\end{notebox}

\subsection{パッケージを追加するには}

\cmd{pyproject.toml} を直接編集してもよいですが、
\cmd{uv add} コマンドを使うと自動で追記されます。

\begin{wslbox}[パッケージの追加]
\begin{lstlisting}[style=bash]
$ uv add pandas
\end{lstlisting}
\end{wslbox}

実行すると \cmd{pyproject.toml} の \cmd{dependencies} に
\cmd{"pandas>=X.X.X"} が自動追記されます。

% -------------------------------------------------
\section{uv.lock とは}
% -------------------------------------------------

\cmd{uv.lock} は uv が自動生成する\textbf{ロックファイル}です。
\textbf{絶対に手動で編集してはいけません。}

\subsection{なぜ必要か}

\cmd{pyproject.toml} には「numpy 1.26 以上」と書きましたが、
これだけでは「今インストールしたら 1.26.4 なのか 1.27.0 なのか」が
わかりません。

時間が経つと最新バージョンが変わるため、
同じ \cmd{pyproject.toml} を使っても
インストールのタイミングによって入るバージョンが変わります。

\begin{center}
\fbox{\parbox{0.85\textwidth}{
\textbf{ロックファイルがないと起きる問題}\\[0.5em]
自分（3月にインストール）：numpy 1.26.4\\
後輩（9月にインストール）：numpy 1.27.2\\
→ 同じコードが異なる環境で動く可能性がある
}}
\end{center}

\subsection{uv.lock が解決すること}

\cmd{uv.lock} には\textbf{全パッケージの正確なバージョンとハッシュ値}が
記録されます。

\begin{wslbox}[uv.lock の中身（抜粋）]
\begin{lstlisting}[style=bash]
version = 1
requires-python = ">=3.11"

[[package]]
name = "numpy"
version = "1.26.4"    # 正確なバージョンが固定されている
source = { registry = "..." }
sdist = { url = "...", hash = "sha256:abc123..." }
\end{lstlisting}
\end{wslbox}

\cmd{uv sync} は \cmd{uv.lock} を読んで\textbf{全員が同じバージョン}を
インストールします。誰がいつ実行しても同じ環境になります。

\subsection{uv.lock の管理ルール}

\begin{itemize}
  \item \textbf{手動で編集しない}（uv が自動管理する）
  \item \textbf{共有フォルダやリポジトリに含める}（全員で共有する）
  \item \cmd{uv add}/\cmd{uv remove} を実行すると自動更新される
  \item \cmd{uv sync} を実行しても \cmd{uv.lock} は変化しない
    （\cmd{uv.lock} に従ってインストールするだけ）
\end{itemize}

\begin{tipbox}[pyproject.toml と uv.lock の関係まとめ]
\begin{tabularx}{\linewidth}{lXX}
  \toprule
  & \cmd{pyproject.toml} & \cmd{uv.lock} \\
  \midrule
  書いてあること & 「これ以上のバージョンが必要」という条件 & 実際にインストールする正確なバージョン \\
  誰が書く & 人間（uv add でも可） & uv が自動生成 \\
  手動編集 & OK & してはいけない \\
  共有する & する & する \\
  \bottomrule
\end{tabularx}
\end{tipbox}

% -------------------------------------------------
\section{uv sync の動作}
% -------------------------------------------------

\cmd{uv sync} は以下のことを\textbf{すべて自動で}行います。

\begin{enumerate}
  \item \cmd{pyproject.toml} を読んで必要な Python バージョンを確認
  \item 必要な Python がなければ自動ダウンロード・インストール
  \item \cmd{.venv/} フォルダを作成（すでにあれば再利用）
  \item \cmd{uv.lock} を読んで記録されたバージョンのパッケージを
    \cmd{.venv/} にインストール
  \item 余分なパッケージ（\cmd{uv.lock} にないもの）があれば削除
\end{enumerate}

\begin{wslbox}[uv sync の出力例（初回）]
\begin{lstlisting}[style=bash]
$ uv sync
Using CPython 3.11.x        <- Python を自動選択
Creating virtual environment at: .venv  <- .venv を作成
Resolved 47 packages in 0.5ms
Installed 47 packages in 8.2s  <- ここが pip より圧倒的に速い
 + bm3d==4.0.2
 + numpy==1.26.4
 + jupyterlab==4.3.6
 ...
\end{lstlisting}
\end{wslbox}

\begin{wslbox}[uv sync の出力例（2回目以降・差分なし）]
\begin{lstlisting}[style=bash]
$ uv sync
Resolved 47 packages in 0.5ms
Audited 47 packages in 1ms   <- 確認のみ、インストールなし
\end{lstlisting}
\end{wslbox}

\begin{notebox}[uv sync は「冪等」]
何度実行しても結果は同じです。
すでに正しい状態なら何もしません。
環境がおかしいと感じたらとりあえず \cmd{uv sync} を試してください。
\end{notebox}
